\par\medskip
\begin{center}
\textls[100]{c) Analyse der Disquisitiones von 1812.}
\end{center}
\noindent\hspace*{\parindent}Neben den eigentlichen Vorarbeiten zu den \textit{Disquisitiones}, über die wir
im Vorhergehenden berichtet haben, befindet sich im Nachlass\footnote{Vergl. auch die Bemerkung von \textsc{Schering}, Werke III, S. 230.} noch eine
in lateinischer Sprache verfasste Handschrift mit dem Titel:
\begin{center}
\textit{Disquisitiones generales}
\end{center}
\noindent\textit{circa functiones transcendentes a serie infinita}
\[
\begin{aligned}
1+\frac{\alpha\mathbin{.}\beta}{1\mathbin{.}\gamma}x
&+\frac{\alpha\mathbin{.}\alpha+1\mathbin{.}\beta\mathbin{.}\beta+1}
{1\mathbin{.}2\mathbin{.}\gamma\mathbin{.}\gamma+1}xx\\
&+\frac{\alpha\mathbin{.}\alpha+1\mathbin{.}\alpha+2\mathbin{.}\beta\mathbin{.}\beta+1\mathbin{.}\beta+2}
{1\mathbin{.}2\mathbin{.}3\mathbin{.}\gamma\mathbin{.}\gamma+1\mathbin{.}\gamma+2}x^3+\text{etc.}
\end{aligned}
\]
\noindent\textit{pendentes}
\begin{center}
auctore \textsc{Carolo Friderico Gauss}\\
Societati regiae traditae Nov. 1811.
\end{center}
\noindent\hspace*{\parindent}Dieses Manuskript hat weder die Bezeichnung »Pars prior« der 1812 gedruckten
Abhandlung, noch die Einteilung in »Sectiones«. Es ist vielmehr
wie folgt gegliedert:
\par
Observationes generales, Art. 1—5.
\par
Relationes inter functiones contiguas, Art. 6—10.\par
\clearpage\SourceLeaf{228}\ArithmeticHead{143}{R}{VIIc. ANALYSE DER DISQUISITIONES.}
\indent Fractiones continuae, Art. 11—13.
\par
De summa seriei [darunter steht durchstrichen: de valore functionis] nostrae
statuendo \(x=1\), Art. 14—27.
\par
Determinatio seriei nostrae per aequationem differentialem secundi ordinis,
Art. 28—42.
\par
Quaedam theoremata specialia, Art. 43—47.
\par
Die Artikel 28—47 dieses Manuskripts bilden die Werke III, S. 205ff.
abgedruckte nachgelassene Abhandlung; die Artikel 1—27 sind von \textsc{Gauss}
umgearbeitet und 1812 als Pars prior veröffentlicht worden. Es entsprechen
die Artikel 1—5 der Handschrift der Introductio, die Artikel 6—10, der ersten,
die Artikel 11—13, der zweiten, die Artikel 14—27, der dritten Sektion der
Veröffentlichung. Namentlich ist die dritte Sektion dadurch erweitert worden,
dass die Theorie der Funktionen \(\Pi\) und \(\Psi\) sehr viel ausführlicher dargestellt
ist als in der Handschrift, entsprechend ist auch die Anzahl der Artikel von
den vierzehn der Handschrift auf dreiundzwanzig vermehrt. Aber beiden Redaktionen
gemein ist die Behandlung der Konvergenzbedingungen für \(F(\alpha,\beta,\gamma,1)\),
von der in den im Nachlass befindlichen Vorarbeiten keine Spur zu finden
ist. \textsc{Gauss} dürfte demnach diese Untersuchung erst im Jahre 1811 durchgeführt
haben. Wenn er am 17. Oktober 1811 an \textsc{Olbers} schreibt (Werke X, 1,
S. 361): »Meine Pallasrechnungen haben nun seit sechs Wochen ganz ruhen
müssen. Ich habe mich viel diese Zeit her mit den transzendenten Funktionen,
worauf die Integration der Gleichung
\begingroup\makeatletter\tagsleft@true\let\veqno\leqno\makeatother
\begin{equation*}
(\alpha+\beta x+\gamma x^2)\frac{d^2y}{dx^2}
+(\delta+\varepsilon x)\frac{dy}{dx}+\zeta y=0\tag{13}
\end{equation*}\endgroup
\noindent führt, beschäftigt und sehr artige Sachen gefunden«, so ist es nicht ausgeschlossen,
dass zu diesen »artigen Sachen« neben der weiteren Ausführung der
Theorie der »inexplikablen Funktionen« \(\Pi\) und \(\Psi\) auch jene Konvergenzuntersuchung
gehört. Jedenfalls war \textsc{Gauss} in der Zeit von Anfang September bis
Mitte Oktober 1811 mit der Ausarbeitung der \textit{Disquisitiones} voll in Anspruch
genommen. Wie aus der Fassung des Titels in dem ursprünglichen Manuskript
hervorgeht, hatte \textsc{Gauss} die Absicht, seine Abhandlung schon im November
1811 der Sozietät der Wissenschaften vorzulegen. Er schreibt am
5. Mai 1812 an \textsc{Bessel} (Werke X, 1, S. 374): »Meine Abhandlung über transzendente
Funktionen habe ich, weil ihr Umfang zu gross wurde, teilen\par
\clearpage\SourceLeaf{229}\ArithmeticHead{144}{L}{L. SCHLESINGER, ÜBER GAUSS' FUNKTIONENTHEORETISCHE ARBEITEN.}
\noindent müssen, \mbox{.\ .\ .}\ Den zweiten Teil meiner Abhandlung hoffe ich auch bald
vollenden zu können\footnote{Vergl. auch die \textit{Anzeige}, Werke III, S. 198 unten.}«. Charakteristisch ist auch die Änderung, die \textsc{Gauss}
im Titel vorgenommen hat; es heisst in der gedruckten Arbeit statt »Untersuchungen
über die von der Reihe \mbox{.\ .\ .\ .}\ abhängenden transzendenten Funktionen«
einfach »Untersuchungen über die Reihe \mbox{.\ .\ .}«.
\par
Man wird wohl annehmen dürfen, dass für die Teilung nicht nur der
äusserliche Grund des allzugrossen Umfangs massgebend war, sondern dass
sowohl die Verzögerung der Mitteilung an die Sozietät um zwei Monate, als
auch der Entschluss, die Arbeit vorläufig gerade an der Stelle abzubrechen,
wo die Differentialgleichung zweiter Ordnung auftreten sollte, durch innere,
sachliche Ursachen bedingt waren. Diese Ursachen können wieder nichts anderes
gewesen sein, als Schwierigkeiten funktionentheoretischer Natur. Wir
werden diesen Punkt weiter unten noch eingehender zu erörtern haben, jetzt
wenden wir uns zur Analyse der 1812 veröffentlichten Abhandlung, an die
wir sogleich auch die der nachgelassenen Fortsetzung anknüpfen wollen. Mit
Rücksicht auf die Erörterungen in dem vorhergehenden Abschnitt, werden wir
bei dieser Analyse die Begriffe Reihe und Konvergenz stets im modernen Sinn
gebrauchen.
\par
Die Konvergenz der Reihe \(F(\alpha,\beta,\gamma,x)\) für beliebige konstante \(\alpha,\beta,\gamma\)
untersucht \textsc{Gauss} (Art. 3) für beliebige komplexe Werte von \(x=a+b\sqrt{-1}\)
durch Vergleichung mit der geometrischen Reihe und findet, dass die Reihe
(sofern \(\gamma\) keine negative ganze Zahl ist) für \(a^2+b^2<1\) stets konvergent, für
\(a^2+b^2>1\) stets divergent ist. Der Fall \(a^2+b^2=1\), speziell \(a=1,\ b=0\)
wird (im III. Abschnitt) besonders untersucht. Es folgt (Art. 4) die Formel
für die Differentiation der Reihe nach \(x\), dann (Art. 5) die Aufzählung von
23 speziellen Fällen, in denen die Reihe bekannte Funktionen darstellt, und
(Art. 6) die Darstellung der Koeffizienten der Entwicklung von
\[
(a^2+b^2-2ab\cos\varphi)^{-n}
\]
nach den Kosinus der Vielfachen von \(\varphi\) durch die Reihe \(F\). Auf diese
Darstellung bezieht sich (vergl. oben Abschnitt V b), S. 90) die Bemerkung
der abgebrochenen Abhandlung über das agM. vom Jahre 1800, Werke III,
S. 371, aus der wir schon oben (S. 91) den Schluss gezogen haben, dass\par
\clearpage\SourceLeaf{230}\ArithmeticHead{145}{R}{VIIc. ANALYSE DER DISQUISITIONES.}
\noindent \textsc{Gauss} jene Abhandlung darum abgebrochen hat, weil ihm\footnote{Vielleicht nachdem er die 1801 erschienene Abhandlung \textsc{Eulers} aus Bd. XII der Nova Acta Petrop. (siehe oben S. 136) kennen gelernt hatte.} der Ausgangspunkt
zu speziell erschien, \textls[100]{und dass wir in den \textit{Disquisitiones} den
ersten Teil des geplanten grossen Werkes über die transzendenten
Funktionen vor uns haben, das auch die ganze Theorie der elliptischen
Funktionen in sich begreifen sollte}. Diese Auffassung wird
durch die folgenden Äusserungen von \textsc{Gauss} bestätigt:
\par
1. In der \textit{Anzeige} der \textit{Disquisitiones} (Werke III, S. 197) heisst es: »Professor
\textsc{Gauss} hat sich mit Untersuchungen über dergleichen höhere transzendente
Funktionen schon seit vielen Jahren beschäftigt \mbox{.\ .\ .}\ Einen verhältnismässig
freilich nur sehr kleinen Teil derselben, der gleichsam als Einleitung
einer künftig zu liefernden Reihe von Abhandlungen angesehen werden
kann, hat er am 30. Januar unter der Aufschrift \textit{Disquisitiones} etc. \mbox{.\ .\ .}\ der
Königl. Gesellschaft der Wissenschaften übergeben«.
\par
2. Die Briefstelle (\textsc{Gauss} an \textsc{Schumacher}, April 1816, Werke X, 1, S. 248):
»In dem zweiten Teile der Abhandlung \textit{Disquisitiones generales circa seriem
infinitam} \(1+\frac{\alpha\mathbin{.}\beta}{1\mathbin{.}\gamma}x\) etc. (welche ich vielleicht bald gebe) werde ich einen Teil
meiner Untersuchungen über die arithmetisch-geometrischen Mittel bekannt
zu machen anfangen«\footnote{Vergl. oben S. 139, 140 den Bericht über den Entwurf aus Handbuch 19, Be, wo übrigens in dem Artikel 6 (Werke X, 1, S. 347) das agM. ausdrücklich genannt wird.}.
\par
Den Grund, weshalb \textsc{Gauss} mit der Reihe \(F(\alpha,\beta,\gamma,x)\) den Anfang macht,
gibt er in der \textit{Anzeige} (Werke III, S. 199) selbst mit folgenden Worten an:
»Allein eben diese erste Erzeugungsart [der Funktion durch die Reihe] führt
schon zu einer Menge merkwürdiger Wahrheiten auf einem bequemen und
gleichsam mehr elementarischen Wege, und deswegen hat der Verf. damit
den Anfang gemacht«. Man kann sich hiernach den Plan, der \textsc{Gauss} vorgeschwebt
hat, einigermassen rekonstruieren. Die Theorie des agM. sollte gleichsam
den Mittelpunkt bilden, indem sie einmal als spezieller Fall der an den
Anfang gestellten Theorie der Reihe \(F(\alpha,\beta,\gamma,x)\) erscheint, während andererseits
von dem agM., d. h. also von der Theorie der Modulfunktion aus, wieder
nach der Seite der doppeltperiodischen Funktionen hin generalisiert werden
sollte, deren Theorie sonach den Abschluss gebildet hätte.\par
\HistoricalFoot{X 2 Abh. 2.}{19}

